\section{Results and Analysis}
\label{sec:results}

Our study has two layers. We first recap, as \emph{background}, the conference-stage influence
analysis that motivates this work: influence in debate is concentrated, stable, and persona-driven
(\S\ref{sec:rq1}); it tracks \emph{stubbornness} rather than accuracy
(\S\ref{sec:rq2}--\S\ref{sec:rq3}); and as a result debate frequently propagates errors
(\S\ref{sec:rq4}), which an influence-aware vote can at best neutralise (\S\ref{sec:rq5}). This layer
is diagnostic and post-hoc; it is not the paper's contribution and we do not build on the
influence-weighted vote.

We then present the paper's five findings. \textbf{RQ1}: multi-agent debate does not merely propagate
errors, it \emph{forgets} correct answers it once held (\S\ref{sec:extinction}). \textbf{RQ2}:
\emph{which} debates fail is predictable, but \emph{which} would benefit from intervention is not, so
neither influence nor failure prediction tells us what action helps (\S\ref{sec:rq2main},
\S\ref{sec:rq13}). \textbf{RQ3}: preserving every hypothesis (never destructively overwriting one)
exposes a large, cheap headroom (\S\ref{sec:preserve}). \textbf{RQ4}: a hypothesis's \emph{provenance}
predicts its correctness (\S\ref{sec:provsel}). \textbf{RQ5}: a zero-cost provenance-aware selector
improves the final decision with no extra inference (\S\ref{sec:provsel}). The hazard-guided
controller that the conference framing would suggest is demoted to a supporting result
(App.~\ref{app:controller}), because it cannot target benefit (RQ2). Unless noted, numbers are pooled
over the five benchmarks $\times$ five seeds under the robust scorer, with cluster-bootstrap 95\% CIs
over questions; in the background layer, $\pm$ denotes the standard deviation across seeds.

\subsection{Influence is concentrated, stable, and persona-driven}
\label{sec:rq1}

Table~\ref{tab:infacc} reports the PageRank influence score and accuracy of each
agent on every benchmark. Two regularities hold across all five tasks. First,
influence is unevenly distributed: the \textit{skeptic} agent is \emph{always} the
least influential participant (influence between $0.144$ and $0.264$), while
\textit{careful} and \textit{bigpicture} consistently occupy the top two positions.
Second, this ordering is remarkably stable across tasks: the average pairwise
Spearman correlation between per-agent influence rankings on different benchmarks is
$0.80$ (minimum $0.50$). Because the three agents differ only in persona and
backbone model, this indicates that influence is largely a property of the
\emph{agent} rather than of the task --- some personas simply move the group more,
regardless of the problem. Figure~\ref{fig:infacc} visualizes the per-agent
influence and accuracy side by side.

\begin{table}[t]
\centering
\small
\caption{Influence score and accuracy per agent (mean$\pm$std over 5 seeds).
\textit{skeptic} is the least influential agent on every benchmark. Role is the
good/bad/low classification from \S\ref{sec:results}.}
\label{tab:infacc}
\begin{tabular}{llccc}
\toprule
Benchmark & Agent & Influence & Accuracy & Role \\
\midrule
\multirow{3}{*}{GSM8K}
 & careful    & $0.376{\pm}.012$ & $0.712{\pm}.017$ & good \\
 & skeptic    & $0.245{\pm}.026$ & $0.602{\pm}.026$ & low  \\
 & bigpicture & $0.379{\pm}.019$ & $0.552{\pm}.021$ & \textbf{bad} \\
\midrule
\multirow{3}{*}{StrategyQA}
 & careful    & $0.426{\pm}.013$ & $0.596{\pm}.030$ & good \\
 & skeptic    & $0.144{\pm}.024$ & $0.574{\pm}.024$ & low  \\
 & bigpicture & $0.430{\pm}.014$ & $0.610{\pm}.054$ & good \\
\midrule
\multirow{3}{*}{CommonsenseQA}
 & careful    & $0.405{\pm}.012$ & $0.592{\pm}.053$ & good \\
 & skeptic    & $0.200{\pm}.011$ & $0.538{\pm}.040$ & low  \\
 & bigpicture & $0.395{\pm}.014$ & $0.548{\pm}.028$ & good \\
\midrule
\multirow{3}{*}{BBH}
 & careful    & $0.335{\pm}.029$ & $0.398{\pm}.033$ & good \\
 & skeptic    & $0.262{\pm}.022$ & $0.314{\pm}.031$ & low  \\
 & bigpicture & $0.404{\pm}.012$ & $0.380{\pm}.048$ & good \\
\midrule
\multirow{3}{*}{TruthfulQA}
 & careful    & $0.347{\pm}.023$ & $0.154{\pm}.039$ & good \\
 & skeptic    & $0.264{\pm}.026$ & $0.072{\pm}.012$ & low  \\
 & bigpicture & $0.389{\pm}.014$ & $0.184{\pm}.021$ & good \\
\bottomrule
\end{tabular}
\end{table}

\begin{figure}[t]
\centering
\includegraphics[width=\linewidth]{results/figures/all_influence_vs_accuracy.png}
\caption{Influence (blue) vs.\ accuracy (orange) per agent on each benchmark. The
influence hierarchy (skeptic lowest) is stable, but its alignment with accuracy is
not.}
\label{fig:infacc}
\end{figure}

\subsection{Influence is not accuracy}
\label{sec:rq2}

A natural hope is that the most influential agent is also the most reliable one, so
that debate amplifies the right voice. Our data does not support this in general.
Across agents, the Pearson correlation between influence and accuracy varies
enormously by task: it is only $0.19$ on GSM8K and $0.24$ on StrategyQA, rising to
$0.74$ on BBH and $0.80$ on TruthfulQA (Fig.~\ref{fig:scatter}). GSM8K is the
starkest case: \textit{bigpicture} is (tied for) the most influential agent
($0.379$) yet the \emph{least} accurate ($0.552$) --- a textbook \emph{bad
influencer} that pulls the group toward its own weaker answers. The same
\textit{bigpicture} agent is a \emph{good} influencer on the other four benchmarks,
where its influence and accuracy are both high. Thus an agent's role is not fixed by
its model but emerges from the model$\times$task interaction, and blindly trusting
influence is only safe when influence and accuracy happen to align.

\begin{figure}[t]
\centering
\includegraphics[width=0.85\linewidth]{results/figures/pooled_influence_vs_accuracy.png}
\caption{Influence vs.\ accuracy pooled over agents and benchmarks. Points are far
from a clean diagonal; on several tasks the most influential agent is not the most
accurate.}
\label{fig:scatter}
\end{figure}

\subsection{Influence comes from stubbornness, not from being right}
\label{sec:rq3}

If influence is not explained by accuracy, what \emph{does} explain it? We test two
candidate drivers over the pooled set of fifteen (agent, benchmark) points: an
agent's \emph{standalone competence}, measured by its round-0 independent accuracy,
and its \emph{stubbornness}, measured as the fraction of opportunities on which it
did \emph{not} change its answer. The contrast is decisive (Table~\ref{tab:source}):
influence is strongly and significantly correlated with stubbornness
($r{=}0.69$, $p{=}0.005$) but essentially uncorrelated with independent accuracy
($r{=}0.15$, $p{=}0.60$). In other words, an agent earns influence in debate by
refusing to move, not by being correct. This is a direct, quantitative
manifestation of the sycophancy and conformity dynamics reported for LLM groups
\cite{b10, b11}: the group gravitates toward whichever agent holds its ground,
independent of whether that ground is true.

\begin{table}[t]
\centering
\small
\caption{What drives influence? Pearson correlation of influence with two candidate
factors, pooled over all (agent, benchmark) pairs. Influence tracks stubbornness,
not standalone competence.}
\label{tab:source}
\begin{tabular}{lcc}
\toprule
Factor & Pearson $r$ & $p$-value \\
\midrule
Independent (round-0) accuracy & $0.15$ & $0.60$ \\
Stubbornness (1 $-$ change rate) & $\mathbf{0.69}$ & $\mathbf{0.005}$ \\
\bottomrule
\end{tabular}
\end{table}

\subsection{Debate frequently propagates errors}
\label{sec:rq4}

Because influence rewards stubbornness rather than correctness, we expect debate to
sometimes destroy answers it should have kept. We quantify this with two
transcript-level events, comparing the round-0 independent answers to the final
majority answer over all $500$ debates per benchmark (Table~\ref{tab:cascade}). A
debate \emph{broke} a question when the independent majority was correct but the
final answer is wrong, and \emph{fixed} it in the opposite case; \emph{net} is
$\text{fixed}-\text{broke}$. The picture is sobering: debate is
\emph{net-negative} on three of five benchmarks, most severely on CommonsenseQA,
where it breaks $107$ questions while fixing only $45$ (net $-62$). We also count
\emph{consensus-to-wrong} events, where at least one agent was correct at round~0
but all agents converge on a wrong final answer --- an outright collapse that occurs
$22$--$89$ times per benchmark. These measurements make concrete the warning of
\cite{b12} that debate need not improve over independent answering, and our
influence analysis explains the mechanism: stubborn-but-wrong agents win the
persuasion process.

\begin{table}[t]
\centering
\small
\caption{Error-cascade accounting over $500$ debates per benchmark. \emph{broke}:
independent majority correct but final wrong; \emph{fixed}: the reverse; \emph{c2w}:
consensus-to-wrong. Debate is net-negative on 3/5 tasks.}
\label{tab:cascade}
\begin{tabular}{lccccc}
\toprule
Benchmark & Fixed & Broke & Net & c2w \\
\midrule
GSM8K         & $56$ & $72$  & $-16$ & $44$ \\
StrategyQA    & $41$ & $29$  & $+12$ & $55$ \\
CommonsenseQA & $45$ & $107$ & $-62$ & $62$ \\
BBH           & $69$ & $86$  & $-17$ & $89$ \\
TruthfulQA    & $24$ & $19$  & $+5$  & $22$ \\
\bottomrule
\end{tabular}
\end{table}

\subsection{Influence-aware voting neutralizes bad influence}
\label{sec:rq5}

Finally, we ask whether the good/bad influence diagnostic is actionable at
aggregation time. Table~\ref{tab:voting} compares majority voting, raw
influence-weighted voting, and accuracy-calibrated influence-aware voting
(Fig.~\ref{fig:voting}). Two effects stand out. First, \emph{naive} influence
weighting is dangerous: because it amplifies whichever agent is most influential ---
which, as \S\ref{sec:rq3} shows, is the most stubborn, not the most accurate --- it
\emph{underperforms} majority voting on GSM8K by $7.6$ points ($0.708\!\to\!0.632$),
exactly the benchmark where the top influencer is a bad influencer. Second,
recalibrating each agent's weight by its role reverses this damage: influence-aware
voting recovers to $0.712$ on GSM8K and matches or slightly exceeds majority voting
on four of five benchmarks (all within one standard deviation), the single exception
being BBH ($-2.4$ points). We therefore do \emph{not} claim influence-aware voting
as a uniform accuracy improvement; rather, it demonstrates that the influence
diagnostic carries actionable signal --- down-weighting agents flagged as bad
influence is enough to convert a harmful aggregation rule into one that is at least
as good as majority voting.

\begin{table}[t]
\centering
\small
\caption{Voting accuracy (mean$\pm$std over 5 seeds). Naive influence weighting hurts
where a bad influencer dominates (GSM8K); accuracy-calibrated influence-aware voting
repairs it and matches majority on 4/5 benchmarks.}
\label{tab:voting}
\begin{tabular}{lccc}
\toprule
Benchmark & Majority & Influence-based & Influence-aware \\
\midrule
GSM8K         & $0.708{\pm}.017$ & $0.632{\pm}.057$ & $\mathbf{0.712}{\pm}.017$ \\
StrategyQA    & $0.592{\pm}.028$ & $0.606{\pm}.045$ & $\mathbf{0.612}{\pm}.040$ \\
CommonsenseQA & $0.596{\pm}.051$ & $0.580{\pm}.052$ & $\mathbf{0.598}{\pm}.053$ \\
BBH           & $\mathbf{0.410}{\pm}.026$ & $0.390{\pm}.039$ & $0.386{\pm}.040$ \\
TruthfulQA    & $0.152{\pm}.034$ & $\mathbf{0.174}{\pm}.033$ & $\mathbf{0.174}{\pm}.033$ \\
\bottomrule
\end{tabular}
\end{table}

\begin{figure}[t]
\centering
\includegraphics[width=\linewidth]{results/figures/all_voting.png}
\caption{Voting strategies per benchmark (mean$\pm$std, 5 seeds). Influence-aware
voting tracks or exceeds majority except on BBH, while naive influence-based voting
collapses on GSM8K.}
\label{fig:voting}
\end{figure}

\paragraph{Summary.} Influence in multi-agent debate is real, stable, and
persona-driven, but it is earned by stubbornness rather than correctness; as a
result debate frequently propagates errors, and only after calibrating influence by
accuracy does an influence-weighted vote stop hurting. Measuring \emph{who}
influences \emph{whom}, and whether that influence is trustworthy, is thus both an
interpretability tool and a safeguard against error cascades.

% ===================================================================
% JOURNAL CONTRIBUTIONS (RQ1-RQ5). The conference influence layer above is
% BACKGROUND. The hazard-guided controller (E6-E12) is demoted to the appendix
% (paper/appendix_controller.tex); a one-paragraph pointer stays in RQ2.
% ===================================================================

\subsection{RQ1 --- Debate forgets correct answers}
\label{sec:extinction}

The background layer showed debate \emph{can} propagate errors; we now show it does something
sharper --- it \emph{forgets answers it already found}. Using a robust extractor (which recovers
\texttt{\textbackslash boxed\{\}} answers hidden by 512-token truncation and \emph{censors}
genuinely-truncated ones so they cannot inflate the count), we measure the accuracy that is
\emph{online-recoverable}: debates in which a correct answer is present at some round but gone from the
final vote. Table~\ref{tab:headroom} shows this behavioural headroom is $+11.6$pp pooled --- roughly
one debate in nine holds the right answer and then discards it, rising to $+15.8$pp on BBH.

\begin{table}[t]
\centering\small
\caption{Recovery headroom (\texttt{e1e4a\_clean}). \emph{Online-recoverable}: a correct answer is
present at some round but not in the final majority; \emph{behavioural}: the same after censoring
genuinely-truncated generations. Debate forgets $11.6$pp of pooled accuracy.}
\label{tab:headroom}
\begin{tabular}{lccc}
\toprule
Benchmark & Final majority & Online-recoverable & Behavioural \\
\midrule
GSM8K         & $73.2\%$ & $8.8\%$  & $7.4\%$  \\
StrategyQA    & $60.0\%$ & $11.0\%$ & $11.0\%$ \\
CommonsenseQA & $63.8\%$ & $12.4\%$ & $12.0\%$ \\
BBH           & $45.6\%$ & $17.8\%$ & $15.8\%$ \\
\midrule
\textbf{Pooled} & $\mathbf{60.7\%}$ & $\mathbf{12.5\%}$ & $\mathbf{11.6\%}$ \\
\bottomrule
\end{tabular}
\end{table}

The loss has a reproducible anatomy. Two behavioural mechanisms drive it: \emph{social adoption} (an
agent copies a peer's wrong answer) and \emph{self-error drift} (an agent revises its own correct
answer away). Social adoption dominates overall ($\approx\!72\%$ of extinctions), but on the
open-answer GSM8K anchor the loss is self-error-led ($57\%$ vs $43\%$), and $73\%$ of adoptions are
\emph{surface}: the answer copies a peer while the reasoning does \emph{not} converge. A lone correct
minority is extinguished $78\%$ of the time. The problem is therefore not a shortage of correct
answers but their destruction --- which is what the rest of the paper attacks.

\subsection{RQ2 --- Predicting failure is not knowing what action helps}
\label{sec:rq2main}

If debate forgets correct answers, the natural fix is to detect a failing debate and intervene. The
detection half is within reach. Global influence \emph{cannot} attribute the mechanism (social vs
self; social-vs-self AUC $0.49\approx$ chance), but a round-0 hazard
detector \emph{predicts} extinction (ROC $0.757$, leave-one-benchmark-out $0.67$--$0.75$), and a
mechanism-agnostic re-query \emph{can} rescue a failing debate (the \texttt{combined} safeguard lifts
preservation $+16$pp pooled at matched call count, gold-free). The targeting half is not. A
hazard-guided controller built on exactly this detector matches always-on intervention at
$\sim\!55\%$ fewer interventions but its edge over a \emph{random} intervention budget is
$\approx\!0$ (its one significant effect is a rescue-$\gg$-harm asymmetry); it \emph{ranks} failures
without converting that ranking into an accuracy win (full battery in App.~\ref{app:controller}). The
next result makes the reason exact.

% E13 — Benefit-Aware controller / uplift-CATE. The paper's sharpest open-problem
% statement. ZERO-GPU on the per-seed E6 cache. Fig: results/figures/e13_benefit_aware.png
% Table: results/tables/e13_benefit_aware.json.
% HONESTY: CATE does NOT beat random (significantly negative on gsm8k); the finding is
% that the treatment effect is UNPREDICTABLE from round-0 features despite a high oracle
% ceiling. Do not spin this as a working controller.

\subsection{Failure prediction is not benefit prediction}
\label{sec:rq13}

The hazard-guided controller (\S\ref{sec:rq2main}; full results in App.~\ref{app:controller}) targets
the E3a \emph{hazard} --- the probability that a debate fails. But what a controller should target is
the \emph{benefit} of intervening, and its robustness battery (App.~\ref{app:controller}) already
hinted the two diverge: the hazard ranks failures, yet targeting it does not beat a random budget. Here we make the distinction
exact. Because the E6 precompute stored \emph{both} counterfactual arms for every debate,
we observe both potential outcomes and therefore the individual treatment effect
$\mathrm{ITE}_i = Y_1^{(i)}-Y_0^{(i)}\in\{-1,0,+1\}$ ($Y_1$: combined-arm recovery, $Y_0$:
control-arm recovery). Any policy intervening on a set $S$ has accuracy
$\mathrm{never}+\tfrac{1}{n}\sum_{i\in S}\mathrm{ITE}_i$, so the optimal policy targets the
largest ITE, and targeting the \emph{true} ITE is an oracle ceiling. We learn the CATE
$\hat\tau(x)=\mathbb{E}[\mathrm{ITE}\mid X{=}x]$ from round-0 features (both a direct ITE
regressor and a two-model T-learner, group-cross-fitted over questions) and compare
policies at matched budget (Figure~\ref{fig:e13}).

Two facts stand out. \emph{First, the ceiling is high}: an oracle targeting true ITE reaches
always-on accuracy at a $\sim\!20\%$ budget, a gain over never of $+16.5$~points pooled
($+11.8$ on GSM8K, $+20.4$ on BBH) --- there is a great deal of benefit to capture.
\emph{Second, none of it is reachable from round-0 features.} The control outcome $Y_0$ is
predictable (OOF AUC $0.73$ --- this is essentially the hazard) and $Y_1$ is predictable
(AUC $0.67$), but their \emph{difference} is not: predicting whether intervention has
\emph{any} effect is barely above chance (AUC $0.58$) and predicting its \emph{sign} fails
outright --- the learned CATE ranks anti-correlated with the true ITE (Spearman $\le 0$ for
both learners, all scopes). Consequently the Benefit-Aware controller does \emph{not} beat a
random budget: pooled $\mathrm{AUC}_{20\text{-}40}$ of CATE$-$random is $-0.003$
(CI $[-0.023,+0.017]$), and on GSM8K it is significantly \emph{negative} ($-0.033$,
CI $[-0.060,-0.006]$), where the anti-correlation actively misleads. It captures
essentially $0\%$ of the oracle gain.

This is the paper's sharpest open-problem statement: \emph{failure prediction and
intervention-benefit prediction are different problems, and only the first is solved from
cheap transcript features.} Differencing two individually-predictable outcomes leaves a
treatment effect that the round-0 behavioral state does not determine --- which is exactly
why the hazard-guided controller matches, but cannot beat, a random budget (App.~\ref{app:controller}).
The large oracle ceiling shows the payoff is real; realizing it needs a benefit signal the
pre-intervention transcript does not carry --- semantic content, mid-debate dynamics, or a
cheap probe of the intervention itself. We therefore position the Benefit-Aware controller
as the correct formulation and its feature side as the central open problem, rather than
overselling a controller that only matches always-on.

\paragraph{Two pre-registered rescue attempts fail (E14, E15).}
One could object that the failure is an artifact of cheap handcrafted features, or of
using only \emph{passive} information. We ran two pre-registered rescues, each with a gate
fixed in advance (predictive: $\mathrm{AUC}_{\text{effect}}>0.65$ or
$\mathrm{Spearman}(\hat\tau,\mathrm{ITE})>0.20$; decision: a cluster-bootstrap CI on
benefit$-$random excluding zero in the $20$--$40\%$ band), same ITE, same learners, no
model search. \textbf{E14} changes \emph{only} the feature vector to residualized
\textsc{nomic-embed} semantic features of each agent's round-0 reasoning (own content, peer
similarity, own-vs-peer gap; per-question centroid removed to kill topic leakage). It fails
far from the gate: effect prediction at chance (AUC $0.47$--$0.50$),
$\mathrm{Spearman}\le 0.06$, benefit$-$random $\mathrm{AUC}_{20\text{-}40}=-0.015$, CI
$[-0.035,+0.004]$ pooled. \textbf{E15} goes further and \emph{actively acquires new
information}: before any revision, each agent is asked to name one concrete flaw in its own
answer and rate its confidence --- a genuine meta-cognitive probe, not a re-encoding of the
transcript. This nudges effect-predictability up (AUC $0.47\!\to\!0.55$ pooled, $0.61$ on
GSM8K) --- active information \emph{is} marginally more informative --- yet still falls far
short of the gate, the sign remains wrong ($\mathrm{Spearman}\le 0$), and the probe
controller does not beat a random budget ($\mathrm{AUC}_{20\text{-}40}=-0.018$, CI
$[-0.038,+0.002]$ pooled; significantly negative on BBH). Across three independent,
pre-registered feature families --- cheap geometry, richer semantics, and actively probed
self-diagnosis --- the treatment effect is not identifiable from the pre-intervention state,
despite the large oracle ceiling. Per the pre-registered stop rule we close the method here
and report the gap itself as the contribution; whether a fundamentally richer
intervention-aware signal (a multi-turn probe, or the intervention's own early trace) can
close it is left to future work.

% ---------------------------------------------------------------------------
\begin{figure}[t]
  \centering
  \includegraphics[width=\linewidth]{results/figures/e13_benefit_aware.png}
  \caption{\textbf{E13 --- benefit prediction is the bottleneck, not failure prediction.}
  Policy value (debate accuracy) vs.\ intervention budget for five targeting rules, pooled
  and per benchmark. An \emph{oracle} that targets the true treatment effect (gold, dashed)
  reaches the ceiling by a $\sim\!20\%$ budget ($+16$--$20$ points over never), so there is
  large benefit to capture. But a \emph{learned} CATE controller (blue) --- trained to
  predict which debates the safeguard \emph{helps} --- tracks at or below random: the
  treatment effect is unpredictable from round-0 features ($Y_0$ and $Y_1$ are each
  predictable at AUC $0.73/0.67$, but their difference is not, AUC $0.58$). Failure
  prediction $\neq$ benefit prediction.}
  \label{fig:e13}
\end{figure}

% ============================================================================
% C3 (preservation ceiling) + C4 (provenance-aware selection).
% \input this after the failure!=benefit section (E13). Numbers: 5 bench x 5 seed,
% gold-free selection, robust scorer; cluster-bootstrap 95% CI over questions.
% ============================================================================

\subsection{Debate forgets: a large preservation headroom}
\label{sec:preserve}

Sections~\ref{sec:rq2main}--\ref{sec:rq13} established that we can predict \emph{which} debates fail but
not \emph{which} would benefit from intervention. We now show that the failure mode itself --- a
correct hypothesis that once existed being lost --- can be attacked \emph{structurally}, without any
targeting. We measure how often the correct answer is present among the \emph{final-round} candidates
versus present \emph{anywhere} across the debate:
\begin{equation}
\mathrm{CR}_{\mathrm{final}} = \Pr\!\big(y^\star \in \{x_i^{T}\}_i\big),\quad
\mathrm{CR}_{\mathrm{hist}} = \Pr\!\big(y^\star \in \textstyle\bigcup_{r\le T}\{x_i^{r}\}_i\big),\quad
\Delta\mathrm{CR}=\mathrm{CR}_{\mathrm{hist}}-\mathrm{CR}_{\mathrm{final}} .
\end{equation}
$\Delta\mathrm{CR}$ is exactly the accuracy a perfect selector over a \emph{preserved} pool would
recover over a perfect selector restricted to the final round. Table~\ref{tab:preserve} shows it is
large and significant on every benchmark: pooled, the correct answer is present in $76.0\%$ of debates
but survives to the final round in only $63.6\%$ --- \textbf{$12.5$ points of correct answers are
generated and then forgotten}, and the total gap to the deployed majority ($50.9\%$) is $+25$pp. The
history pool contains only $\approx 3$ distinct hypotheses per question, so the recovered candidates
are not noise. The effect is largest on BBH ($+19.6$pp) --- precisely the benchmark where the
hazard-guided controller failed (App.~\ref{app:controller}) --- showing that preservation helps most exactly
where targeting helps least.

\begin{table}[t]
\centering\small
\caption{Preservation ceiling. Correct-answer recall at the final round vs.\ across the full history,
and the deployed majority accuracy. $\Delta\mathrm{CR}$ is the forgotten-correct rate (cluster-boot
95\% CI over questions). Pool = mean distinct hypotheses per question.}
\label{tab:preserve}
\begin{tabular}{lccccc}
\toprule
Benchmark & Maj.\ acc & $\mathrm{CR}_{\mathrm{final}}$ & $\mathrm{CR}_{\mathrm{hist}}$ & $\Delta\mathrm{CR}$ [95\% CI] & Pool \\
\midrule
GSM8K         & 0.732 & 0.824 & 0.920 & $+0.096\,[.070,.128]$ & 2.40 \\
BBH           & 0.456 & 0.602 & 0.798 & $+0.196\,[.156,.236]$ & 2.84 \\
StrategyQA    & 0.600 & 0.750 & 0.908 & $+0.158\,[.112,.206]$ & 1.64 \\
CommonsenseQA & 0.638 & 0.806 & 0.936 & $+0.130\,[.100,.162]$ & 2.26 \\
TruthfulQA    & 0.120 & 0.196 & 0.240 & $+0.044\,[.022,.068]$ & 6.31 \\
\midrule
\textbf{Pooled} & \textbf{0.509} & \textbf{0.636} & \textbf{0.760} & $\mathbf{+0.125\,[.108,.142]}$ & 3.09 \\
\bottomrule
\end{tabular}
\end{table}

\subsection{Provenance recovers part of the headroom at zero extra cost}
\label{sec:provsel}

Realizing the ceiling requires \emph{selecting} the correct hypothesis from the preserved pool without
gold. We first rule out two tempting selectors. A deterministic arithmetic verifier (accept a candidate
whose value equals a numerically valid expression in its own reasoning) \emph{hurts} versus simple
frequency ($-0.032$ on GSM8K): correct and incorrect answers are equally arithmetically self-consistent
because errors are in problem \emph{formulation}, not computation. An external LLM judge is worse: with
the judge forced to pick a pool label (isolating preservation as
$\Delta_{\mathrm{HP}}=\mathrm{Acc}[J(\mathcal{H}_{\mathrm{hist}})]-\mathrm{Acc}[J(\mathcal{H}_{\mathrm{final}})]$),
both a 7B and a 14B judge fail the gate ($\Delta_{\mathrm{HP}}=-0.047$ and $-0.018$); tellingly, the
14B judge solving the question directly (0.76) far outscores it selecting from the pool (0.53), so a
capable external judge is a \emph{liability} under capability mismatch, not a fix.

The signal that does work is intrinsic: a hypothesis's own \emph{provenance}. Table~\ref{tab:prov}
shows that whether an answer was produced by many agents \emph{independently} (before seeing peers),
how long it \emph{persisted}, and when it \emph{first appeared} each predict its correctness at the
candidate level, monotonically and across benchmarks. An answer produced independently by three agents
is $83\%$ correct; one that only \emph{emerges} during the debate is $\approx 7\%$ correct --- the
majority vote discards this distinction entirely.

\begin{table}[t]
\centering\small
\caption{Provenance predicts correctness (candidate level; directional AUC with cluster-boot CI, and
$\Pr(\text{correct})$ by value). Independent support and persistence are monotone across benchmarks.}
\label{tab:prov}
\begin{tabular}{lcccl}
\toprule
Signal & dir.\ AUC [95\% CI] & mono.\ & $\Pr(\text{correct}\mid\cdot)$ \\
\midrule
$N_{\mathrm{indep}}$ (round-0 support) & $0.823\,[.800,.845]$ & 4/5 & $0{:}.07\;\;1{:}.17\;\;2{:}.58\;\;3{:}.83$ \\
$N_{\mathrm{persist}}$ (rounds alive)  & $0.806\,[.785,.824]$ & 5/5 & $1{:}.06\;\;2{:}.30\;\;3{:}.57$ \\
$N_{\mathrm{support}}$ (total votes)   & $0.857\,[.836,.877]$ & 2/5 & $.05 \rightarrow .84$ \\
first round (late emergence)           & $0.695\,[.681,.709]$ & 4/5 & $r_0{:}.37\;\;r_1{:}.08\;\;r_2{:}.04$ \\
\bottomrule
\end{tabular}
\end{table}

Turning this into a selector that picks one hypothesis per debate,
$S(h)=w_1 N_{\mathrm{indep}}+w_2 N_{\mathrm{persist}}+w_3 N_{\mathrm{support}}-w_4\,\mathrm{first\_round}$
(equal weights, or logistic weights cross-fit by question so no weight ever sees its own test fold),
gives the ladder in Table~\ref{tab:selector}. The provenance selector beats the deployed final-majority
by $+3.2$pp pooled (up to $+8$pp) and beats the conference influence-weighted vote --- which itself
\emph{loses} to plain majority ($-2.6$pp) --- all with \textbf{no additional inference}. It ties the
round-0 independent majority, because trusting the pre-debate answer is already a strong rule: debate
is \emph{net-destructive} on four of five benchmarks. A convergence-gated variant, which credits late
consensus only when the supporters' reasoning genuinely converges (lexical similarity residualized
against the debate's own baseline, controlling for topic), correctly recovers the one
constructive-debate benchmark (StrategyQA, $0.580\!\to\!0.604$) but cannot beat round-0 pooled --- a
single global rule cannot serve both regimes, and the residual requires per-instance knowledge of
whether a revision is constructive, which is precisely the mechanism-attribution problem
(\S\ref{sec:rq2main}) that is at chance. The remaining $\approx 86\%$ of the ceiling is our stated open
problem.

\begin{table}[t]
\centering\small
\caption{Inference-time selection rules over the frozen debate logs (gold-free; cluster-boot CI over
questions). Provenance beats the deployed majority and the conference vote at zero extra LLM calls;
an external judge and influence-voting lose. Oracle = a perfect selector over the preserved pool.}
\label{tab:selector}
\begin{tabular}{lcccl}
\toprule
Selector & Pooled acc & $\Delta$ vs.\ majority [95\% CI] & Extra calls \\
\midrule
Final majority (deployed)          & 0.528 & --- & 0 \\
Conference influence-voting        & 0.502 & $-0.026\,[-.036,-.018]$ & 0 \\
Round-0 (independent) majority     & 0.560 & $+0.032\,[.015,.048]$ & 0 \\
\textbf{Provenance selector (ours)}& \textbf{0.564} & $\mathbf{+0.035\,[.022,.049]}$ & \textbf{0} \\
\midrule
Oracle (history-pool ceiling)      & 0.760 & $+0.232$ & --- \\
\bottomrule
\end{tabular}
\end{table}


% ===================================================================
% APPENDIX (demoted): full hazard-guided controller battery E6-E9,E12.
% Relocate under \appendix in the master document if desired.
% ===================================================================
% Appendix: full hazard-guided controller battery (E6 controller + E7-E9,E12 robustness).
% DEMOTED supporting material. The controller MATCHES always-on at ~half the interventions but its
% edge over a random budget is ~0 -- which is exactly why RQ2 concludes benefit is not identifiable.
% Kept for completeness/reproducibility; not part of the main arc. Relocate under \appendix in the
% master document if desired.
\section{Hazard-Guided Controller: Full Results}
\label{app:controller}

RQ2 (\S\ref{sec:rq2main},~\S\ref{sec:rq13}) concludes that intervention \emph{benefit} is not
identifiable from the pre-intervention state. This appendix reports the full controller battery behind
that conclusion: a hazard-guided controller that selectively intervenes on the debates its detector
flags. The headline is deliberately modest --- the controller \emph{matches} always-on intervention at
$\sim\!55\%$ fewer interventions, but its edge over a \emph{random} intervention budget is within
noise, and its only significant effect is that rescues outnumber harms. We keep it for completeness and
because its per-seed statistics (E7--E9, E12) establish the cluster-bootstrap robustness protocol
reused throughout the paper.

% E6 — hazard-guided adaptive controller. Slots into results_analysis.tex as the
% closing subsection (\S\ref{sec:rq6}). Figures: results/figures/e6_controller.png,
% results/figures/e6_collateral.png. Source table: results/tables/e6_controller.json.
%
% Terminology discipline (do NOT revert):
%   * "hazard-guided controller" — NOT "influence-guided". E2 showed influence alone
%     is insufficient to attribute the mechanism; at the g=0 gate the detector uses
%     round-0 disagreement geometry (F1) + text (F4) + influence (F3). F2 dynamics are
%     ZERO at g=0 (nothing has revised yet) -- do not claim "revision dynamics" here.
%   * tau* is an OBSERVED Pareto point, not a tuned knob (E12 cross-fits it).
%   * NUMBERS ARE PER-SEED (fixed cache, keyed bench:qid:seed). The earlier per-qid
%     cache collapsed the 5 seeds -> inflated headroom (8.5->4.5pp) and rescue
%     (.46->.40) and a spurious gsm8k advantage. Do NOT reintroduce old numbers.
%   * Honest headline (E9, cluster-bootstrap over 184 questions): the controller
%     MATCHES always-on (hazard-always = +0.0pp) at ~55% fewer interventions; its edge
%     over random is borderline (+2pp pooled, ~0 within each benchmark); the ONE
%     significant result is the rescue>>harm asymmetry (0.40 vs 0.21).

\subsection{Selective intervention is efficient: a hazard-guided controller}
\label{sec:rq6}

The preceding results close a chain of inference. Debate can \emph{destroy} a correct
answer that was alive early (\S\ref{sec:rq4}, E1); this harm has a reproducible
anatomy (E4a); the round-0 dynamics make it \emph{predictable} before it happens
(E3a); and a targeted re-query \emph{repairs} it when it is about to occur (E5). The
open question is whether prediction and repair can be combined into a policy that is
better than intervening everywhere. Crucially, the triggering signal is \emph{not}
influence: E2 established that influence alone does not attribute the failure
mechanism. The controller is therefore \emph{hazard-guided} --- it fires on the E3a
Stage-A hazard, a leave-one-group-out out-of-fold score built from the round-0
behavioral state: answer-disagreement geometry (F1), reasoning-text features (F4), and
global influence (F3). Revision dynamics (F2), which strengthen the pooled detector at
later observation times, are by construction unavailable this early --- the earliest
deaths occur before any agent has revised. \S\ref{sec:rq7} shows that \emph{no single
channel dominates}: structure, text, and influence each contribute, influence is
neither necessary nor sufficient, and the combined detector is the robust choice.

\paragraph{Setup.}
We evaluate on the $n=516$ debates ($184$ distinct questions $\times$ up to five seeds)
that have a correct answer alive at round~0 --- the population that \emph{has} something
to protect (extinction prevalence $22.5\%$). For each debate we precompute, \emph{per
seed}, two counterfactual continuations from a single early gate: \textsc{control} (a
plain stochastic re-query, the E5 baseline) and \textsc{combined} (self-verification
plus an influence-resistant re-query). Every controller is then an \emph{offline}
selection policy over these two stored outcomes, so the entire accuracy--intervention
curve is swept at zero additional GPU cost. Two anchors bound the curve: a \emph{never}
policy (\textsc{control} everywhere; accuracy $58.3\%$) and an \emph{always-on} policy
(\textsc{combined} everywhere; accuracy $62.8\%$), a headroom of $4.5$ points. The
\emph{hazard-guided} policy
applies \textsc{combined} iff the E3a hazard exceeds a threshold $\tau$; a
\emph{random-matched} policy spends the identical intervention budget on a random
subset and is our null model for ``does the detector add anything.''

\paragraph{Selective beats indiscriminate.}
Figure~\ref{fig:e6}A traces accuracy against intervention rate. In point estimate the
hazard-guided controller tracks above the random-matched budget across the sweep and
\emph{numerically exceeds} always-on at a much reduced intervention rate: at its
observed peak operating point ($\tau^\star=0.52$) it reaches
\[
  \mathrm{Acc}_{\text{hazard}} = 63.2\%
  \;\gtrsim\;
  \mathrm{Acc}_{\text{always}} = 62.8\%
  \;>\;
  \mathrm{Acc}_{\text{never}} = 58.3\%,
\]
while intervening on only $45\%$ of debates ($55\%$ fewer safeguard
\emph{interventions} than always-on). Under a cluster-bootstrap over the $184$
independent questions (\S\ref{sec:rq9}), the hazard$-$always margin is $+0.0$ points
(CI $[-1.6,+1.2]$), squarely within noise, so the defensible claim is that the
controller \emph{matches} always-on accuracy at roughly half the interventions; its
edge over a matched-random budget is \emph{suggestive but not conclusive} ($+2$ points
pooled, near zero within either benchmark). The one comparison that \emph{is}
significant is the collateral asymmetry below. This is not a compute
argument: the combined safeguard adds only $\sim\!0.06$ model calls per debate on
average ($5.65\!\to\!5.72$), so raw token cost is essentially unchanged; the
controller's value is doing the job of always-on with far fewer interventions and less
collateral harm. Figure~\ref{fig:e6}B recasts this as
gain retention, $(\mathrm{Acc}-\mathrm{never})/(\mathrm{always}-\mathrm{never})$.
Because the always$-$never headroom is only $4.5$ points, gain retention is a
high-variance summary; we therefore rely on the raw accuracy edge over the
matched-random budget ($\approx\!+2$ points pooled at a $40\%$ intervention rate) with
its cluster-bootstrap interval (\S\ref{sec:rq9}) and on the baseline comparisons
(\S\ref{sec:rq7}), rather than on any single gain-retention number.

We present $\tau^\star$ as an \emph{observed} Pareto operating point, not a tuned
hyperparameter, and \S\ref{sec:rq12} confirms by cross-fitting that a data-chosen
threshold does not inflate the result (held-out accuracy matches always-on). The
conclusion rests on the full budget frontier, not any single threshold.

\paragraph{Why a controller has to exist.}
The gain is not free efficiency --- it reflects that intervention is
\emph{asymmetric} (Figure~\ref{fig:e6collateral}). Applying the combined safeguard
indiscriminately rescues $39.5\%$ of would-die debates ($n=215$, those the plain
re-query leaves wrong) but \emph{harms} $20.6\%$ of would-live debates ($n=301$,
those that were already correct and are knocked off the answer by the extra
perturbation). Always-on intervention pays this collateral tax on every safe debate;
the hazard-guided controller exists precisely to concentrate intervention where its
expected value is positive and to withhold it where the dominant risk is
self-inflicted harm. This asymmetry (rescue $\gg$ harm) is the one result that
\emph{is} statistically significant under question-level clustering (\S\ref{sec:rq9}),
and it is the controller's principled justification: even where selective and
indiscriminate intervention tie on aggregate accuracy, targeting is preferable because
it avoids the collateral damage that blanket safeguarding necessarily incurs.

% ---------------------------------------------------------------------------
\begin{figure}[t]
  \centering
  \includegraphics[width=\linewidth]{results/figures/e6_controller.png}
  \caption{\textbf{A hazard-guided controller makes selective intervention efficient}
  ($n=516$ debates with a correct answer alive at round~0).
  \textbf{(A)} Debate accuracy vs.\ intervention rate. In point estimate the
  hazard-guided controller (blue) tracks above a random-matched budget (grey, within-run
  $95\%$ CI) and \emph{numerically exceeds} always-on: at the observed operating point
  $\tau^\star=0.52$ (star) it reaches $63.2\%$ vs.\ the always-on $62.8\%$, while
  intervening on only $45\%$ of debates (\emph{never} floor $58.3\%$). Under a proper
  cluster-bootstrap over questions (\S\ref{sec:rq9}) the hazard$-$always margin is $+0.0$
  points (within noise) and the hazard$-$random edge is borderline; the robustly
  significant result is the collateral asymmetry (Fig.~\ref{fig:e6collateral}).
  \textbf{(B)} The same sweep as gain retention. In the $20$--$40\%$ operating band
  (shaded) the controller clears the $\mathrm{GR}=0.70$ gate in point estimate and lies
  above the matched-random budget; its vertical gap over the grey curve is the
  detector's (point-estimate) value. $\tau^\star$ is an observed Pareto point, not a
  tuned threshold (cross-fitted in \S\ref{sec:rq12}).}
  \label{fig:e6}
\end{figure}

\begin{figure}[t]
  \centering
  \includegraphics[width=0.72\linewidth]{results/figures/e6_collateral.png}
  \caption{\textbf{Intervention is asymmetric --- the controller's reason to exist.}
  Applying the combined safeguard everywhere rescues $39.5\%$ of would-die debates
  ($n=215$) but inflicts collateral harm on $20.6\%$ of would-live debates
  ($n=301$); the two intervals do not overlap (\S\ref{sec:rq9}), so the asymmetry is
  significant. Because indiscriminate safeguarding taxes every already-correct debate,
  a hazard-guided policy that withholds intervention on low-hazard debates does the job
  of always-on with far fewer interventions.}
  \label{fig:e6collateral}
\end{figure}

% E7-E9, E12 — robustness battery for the hazard-guided controller. All ZERO-GPU,
% built on the frozen E6 precompute cache via experiments/e6_offline.py.
% Figures: results/figures/{e7_baselines,e8_ablation,e9_stats,e12_threshold}.png
% Tables : results/tables/{e7_baselines,e8_ablation,e9_stats,e12_threshold}.json
%
% HONESTY CONTRACT (do not soften back into overclaims):
%   * controller MATCHES always-on (not "beats"); hazard-always is within noise.
%   * detector edge over random is BORDERLINE and gsm8k-driven; bbh shows nothing.
%   * the significant result is rescue >> harm asymmetry.
%   * effective n = 184 questions (seed-collapse in the per-qid cache).

\subsection{Is the controller real? Baselines and ablation}
\label{sec:rq7}

A hazard-guided controller is only worth its complexity if it beats trivial policies
and if its gain is a genuine per-task signal rather than an artifact. We test both on
the same $n=516$ offline population, comparing every policy at a \emph{matched}
intervention budget (apply the safeguard to the top-$r$ debates by risk score, ties
broken at random).

\paragraph{Baselines (E7).}
Figure~\ref{fig:e7} pits the learned hazard against five deployable heuristics. The
round-0 \emph{disagreement} geometry is a strong baseline: intervening where agents
disagree, where answer-entropy is high, or where the majority margin is small all reach
$\approx\!61.3\%$ accuracy at a $40\%$ budget, versus $60.1\%$ for random. The learned
hazard is the \emph{best} single signal --- highest at every budget ($62.2\%$ at $40\%$)
and, more tellingly, the best ranker of true extinctions (AUPRC $0.375$ vs.\ $\le0.33$
for any heuristic and $0.225$ for random) --- but the margin over the structural
heuristics is small in the operating band and widens only at higher budget. Two
negative controls sharpen the story: an \emph{influence-only} detector (the F3
PageRank/stub\-bornness block) \emph{underperforms} the structural heuristics
(AUPRC $0.300$), and a \emph{revision} trigger is not even evaluable at the $g=0$ gate
--- no agent has revised yet --- so it coincides with random. The controller's power is
round-0 behavioral geometry, not influence and not revision dynamics.

\paragraph{Ablation (E8).}
Figure~\ref{fig:e8} decomposes the detector. Pooled (Fig.~\ref{fig:e8}A), no single
feature group dominates and \emph{influence is not required}: dropping the F3 block
actually raises AUPRC ($0.375\!\to\!0.415$), while text-surface features (F4) carry the
most unique pooled signal. But the pooled view is confounded --- GSM8K and BBH differ in
both extinction prevalence and verbosity --- so we validate \emph{within} each benchmark
(Fig.~\ref{fig:e8}B). Here the honest picture emerges: the detector \emph{ranks}
extinctions within GSM8K (the full detector attains AUPRC lift $1.98\times$; influence
alone $2.46\times$), yet \emph{targeting} those debates yields essentially no accuracy
edge over a random budget at a $40\%$ intervention rate --- $+0.0$ points within GSM8K,
$-1.2$ points within BBH, and only $+2.1$ points pooled. (An earlier per-question cache
that shared one outcome across a question's five seeds had inflated the GSM8K figure;
the numbers here use the corrected per-seed outcomes.) Ranking ability does not convert
into a controller advantage because the safeguard's rescue rate is modest and it harms
would-live debates (\S\ref{sec:rq9}). Influence remains benchmark-dependent (best ranker
in GSM8K, $\approx$random in BBH). For contrast, the pooled multi-observation-time E3a
detector, where the F2 dynamics block \emph{is} observable, improves monotonically
(AUPRC $0.19\!\to\!0.22\!\to\!0.24$ as dynamics then influence are added); dynamics help
\emph{when} available, but the controller's early gate cannot see them. The robust
reading: combine all channels for ranking, but the controller currently matches --- does
not beat --- a random budget on accuracy in either benchmark.

\subsection{Statistical robustness}
\label{sec:rq9}

Each debate now carries its own \emph{per-seed} counterfactual outcome (an earlier cache
that shared one outcome across a question's five seeds has been corrected). The five
seeds of a question are still correlated repeats of the same question, so the resampling
unit is the question, not the debate: the $516$ debates form $184$ clusters. We report a
cluster-bootstrap (resampling questions, $B=4000$, the deployed detector held fixed);
Figure~\ref{fig:e9} shows the intervals. Three conclusions follow. (i)~The controller
\emph{matches} always-on: the hazard$-$always accuracy gap at the peak operating point is
$+0.0$~points, CI $[-1.6,+1.2]$ --- squarely within noise. (ii)~The detector's edge over
a matched-random budget is \emph{suggestive but not conclusive}: the area under the
policy curve in the $20$--$40\%$ band is $+0.020$, CI $[-0.003,+0.042]$; every per-budget
interval includes zero, pooled and within each benchmark (no benchmark reaches
significance individually). (iii)~The one robustly significant effect is the
\emph{asymmetry} of intervention: applied everywhere, the safeguard rescues $39.5\%$ of
would-die debates (CI $[32.4,46.8]$) while harming $20.6\%$ of would-live debates
(CI $[15.5,26.1]$) --- non-overlapping intervals (and likewise within GSM8K). The
\emph{principle} of selective intervention is thus on firm ground; the current detector
ranks extinctions but does not yet convert that into a significant accuracy edge over a
random budget --- an honest gap for the mechanism, not the principle.

\subsection{Threshold and budget robustness}
\label{sec:rq12}

Finally we rule out the natural objection that $\tau^\star=0.52$ was chosen on the test
set (Figure~\ref{fig:e12}). Our recommended headline is not a threshold at all but the
\emph{budget-conditioned frontier} (Fig.~\ref{fig:e12}A): ``given an intervention budget
$r$, apply the safeguard to the $r$ highest-hazard debates.'' A budget is fixed a priori,
so nothing is tuned on test, and the frontier itself is the deliverable. To test whether
a \emph{data-chosen} threshold would generalize, we cross-fit it (GroupKFold-$5$ over
questions; on each training split pick $\tau^\star=\arg\max_\tau$ train accuracy and apply
it to the held-out split). The selected threshold is unstable across folds
($\tau^\star\in[0.004,0.52]$) and the pooled held-out accuracy is $61.4\%$,
CI $[56.2,66.7]$, at a mean $67\%$ intervention rate --- below the in-sample peak
($63.2\%$) and indistinguishable from always-on ($62.8\%$). The peak is therefore mildly
selection-optimistic, and a properly cross-validated policy simply matches always-on ---
exactly consistent with \S\ref{sec:rq9}.

\paragraph{Limitations and next steps.}
Three honest limitations bound these claims. \emph{First}, on aggregate accuracy the
controller \emph{matches} both always-on and a matched-random budget --- the edge over
random is not significant on either benchmark. Its demonstrated value is doing the job of
always-on at roughly half the interventions while avoiding a statistically significant
collateral harm; the accuracy payoff of \emph{targeting} (over random targeting) is a
real, open gap. \emph{Second}, only two benchmarks are covered here; a separate
detectability sweep (E11) shows extinction is even more predictable on StrategyQA
(hazard ROC $0.87$) than GSM8K, so a benchmark expansion is the natural next test of
whether targeting ever converts into accuracy. \emph{Third}, the three-agent ensemble is
fixed; whether the phenomenon survives other model families and $N>3$ is untested (E10).
(The per-question seed-collapse present in an earlier draft has been fixed here: every
seed now carries its own counterfactual outcome, which is what shrank the apparent
always-on headroom from $8.5$ to $4.5$ points and removed a spurious GSM8K advantage.) We
regard the mechanism --- predict-then-target, justified by the significant rescue/harm
asymmetry and a detector that ranks extinctions at $\sim\!2\times$ prevalence --- as the
contribution, and converting that ranking into a significant accuracy edge as the main
open problem, which \S\ref{sec:rq13} localizes precisely: the treatment \emph{benefit},
unlike the failure, is not predictable from round-0 features.

% ---------------------------------------------------------------------------
\begin{figure}[t]
  \centering
  \includegraphics[width=0.82\linewidth]{results/figures/e7_baselines.png}
  \caption{\textbf{E7 --- controller baselines at matched budget.} The learned hazard is
  the best single signal at every budget and the best extinction ranker (AUPRC $0.375$),
  but round-0 disagreement/entropy/margin heuristics are strong; \emph{influence-only}
  trails them and a \emph{revision} trigger is degenerate at the $g=0$ gate
  ($\equiv$ random). Bands are within-run random CIs.}
  \label{fig:e7}
\end{figure}

\begin{figure}[t]
  \centering
  \includegraphics[width=\linewidth]{results/figures/e8_ablation.png}
  \caption{\textbf{E8 --- ablation.} \textbf{(A)} Pooled, no group dominates and influence
  (F3) is not required (dropping it raises AUPRC). \textbf{(B)} Within-benchmark accuracy
  edge over a random budget at a $40\%$ intervention rate is near zero everywhere
  ($+0.0$ on GSM8K, $-1.2$ on BBH, $+2.1$ pooled): the detector ranks extinctions but
  targeting them barely beats random. Influence is benchmark-dependent (best on GSM8K,
  $\approx$random on BBH).}
  \label{fig:e8}
\end{figure}

\begin{figure}[t]
  \centering
  \includegraphics[width=\linewidth]{results/figures/e9_stats.png}
  \caption{\textbf{E9 --- cluster-bootstrap over $184$ questions.} \textbf{(A)}
  hazard$-$random is borderline (every interval includes zero, pooled and within each
  benchmark); hazard$-$always is within noise ($+0.0$ points). \textbf{(B)} the
  rescue$\gg$harm asymmetry is significant (non-overlapping intervals, pooled and GSM8K).
  The principle is solid; converting the detector's ranking into an accuracy edge is the
  open gap.}
  \label{fig:e9}
\end{figure}

\begin{figure}[t]
  \centering
  \includegraphics[width=\linewidth]{results/figures/e12_threshold.png}
  \caption{\textbf{E12 --- threshold/budget robustness.} \textbf{(A)} We report the
  budget-conditioned frontier (top-$r$ by hazard), which needs no test-time tuning.
  \textbf{(B)} A cross-fitted, data-chosen threshold yields $61.4\%$ held-out accuracy ---
  below the $63.2\%$ in-sample peak and $\approx$ always-on ($62.8\%$): the peak is mildly
  selection-optimistic, no cherry-pick advantage survives.}
  \label{fig:e12}
\end{figure}


